\section{1991--2007: A Movement Without a Founder, and the Question of
Formal Adherence}

\subsection{Succession inside the Society}

Archbishop Lefebvre died on 25 March 1991; Bishop Bernard Fellay was
elected Superior General in 1994 and held the office until 2018, when the
Society's general chapter elected Father Davide Pagliarani. These are
internal facts of the Society, reported by the Society, and this account
carries them as such; no Roman act establishes them.\endnote{The
Society's own institutional history and its published chapter
communiqués, read 2026-07-25. Note that the Superior General after 2018
is a priest, not a bishop: the Society's bishops hold no office of
government in it. Its own statistical page states that the ministerial
function of its bishops ``being limited to the administration of the
sacraments of Holy Orders and Confirmation, our bishops neither receive
nor claim any episcopal jurisdiction over priests or the faithful''
(sspx statistics for 1 November 2025, published 4 January 2026).}

That last point is more than an internal curiosity. From the beginning
the Society has described its bishops as auxiliaries without
jurisdiction---bishops who ordain and confirm but govern nothing---and
this description recurs verbatim in its statement of 1 July 2026. It is
the Society's own characterization of its structure, not a canonical
status recognized by the Holy See, and this account never treats it as
one.

\subsection{The 1996 note: what ``formal adherence'' means, and who said so}

\work{Ecclesia Dei} 5~c) had said that formal adherence to the schism
carries excommunication under canon 1364. Bishops asked what that meant
in practice. In July 1996 the Bishop of Sion---the Swiss diocese in which
Écône lies---wrote to the Congregation for Bishops asking for an
authoritative interpretation, after confused press reports; the
Congregation asked the Pontifical Council for Legislative Texts for its
opinion, and on 24 August 1996 the Council's president, Archbishop Julián
Herranz, sent a covering letter and an attached note, which the Council
then published.\endnote{Pontifical Council for Legislative Texts, note of
24 August 1996 (Italian text on the Holy See website;
\work{Communicationes} 29 [1997] 239--243), read 2026-07-25. The covering
letter identifies the request of Bishop Norbert Brunner of Sion,
transmitted by letter of 26 July 1996.}

The document's own account of its authority deserves quoting in
substance, because it is regularly cited as though it were a definitive
Roman ruling. The Council's president wrote that the problem raised by
the Ordinary of Sion did not appear to require an authentic
interpretation, either of \work{Ecclesia Dei}, or of the Congregation's
decree of 1 July 1988, or of the relevant canons 1364 \S\,1 and 1382,
because no genuine \latin{dubium iuris} had been shown; nevertheless, to
assist the Congregation, the attached note offered ``some considerations
and suggestions.'' In other words: a consultative opinion of a curial
council, expressly not an authentic interpretation.

Its contents are nonetheless precise, and they became binding-by-adoption
thirty years later. The note states that the schism was declared in
immediate relation to the episcopal ordinations of 30 June 1988, but that
this gravest act of disobedience consummated a progressive global
situation of a schismatic character; that until changes restore the
necessary hierarchical communion, the whole Lefebvrian movement is to be
held schismatic, a formal declaration of the supreme authority existing in
the matter; that the validity of the bishops' excommunications cannot
reasonably be doubted, no exempting or attenuating circumstance appearing;
and that as to the state of necessity Lefebvre believed himself to be in,
such a state must be verified objectively, and there is never a necessity
to ordain bishops against the will of the Roman Pontiff, head of the
college of bishops---which would mean the possibility of ``serving'' the
Church by an attack on her unity.\endnote{Same note, nn.~1--4. Paraphrase
from the Italian; no project translation is offered. The note's
n.~4 also declines to assess a canonist's contrary thesis reported in the
press, because that thesis was not known to the Council.}

On formal adherence the note gives the two-element test that has governed
the question ever since: an internal element, freely and consciously
sharing the substance of the schism by placing that option above obedience
to the pope; and an external element, the externalization of that option,
whose most manifest sign is exclusive participation in Lefebvrian
``ecclesial'' acts without taking part in the acts of the Catholic
Church---though the note immediately adds that this sign is not
univocal, since some may attend such liturgies without sharing the
schismatic spirit. For deacons and priests of the movement, it holds that
ministerial activity within the movement is more than evident proof that
both requirements are met. For other faithful, occasional participation
in liturgical acts or activities, without adopting the movement's
attitude of doctrinal and disciplinary disunion, is not sufficient; the
various situations must be judged case by case, in the competent forums
of external and internal forum. Finally, it insists on distinguishing the
moral question of the sin of schism from the juridical-penal question of
the delict and its sanction, to which the Code's Book VI---including the
canons on exempting and attenuating factors---applies.\endnote{Same note,
nn.~5--8, again paraphrased.}

\actrecord{Explanatory note of the Pontifical Council for Legislative
Texts, 24 August 1996.}{The two-element test for ``formal adherence''; the
Council's view that the whole movement is to be held schismatic; its view
that the 1988 excommunications' validity cannot reasonably be doubted;
and its insistence that lay situations be judged case by case.}{A
consultative note of a curial council, sent to another dicastery, which
expressly denies that an authentic interpretation was required. Until
2026 it bound no one by its own force. The Dicastery for the Doctrine of
the Faith made it its own on 2 July 2026, and that adoption---not the
1996 text---is what gives it force today.}

\subsection{Jubilee contacts and \work{Summorum Pontificum}}

The Society made a public pilgrimage to Rome in the Jubilee year 2000,
and contacts resumed. The next Roman act that changed the landscape was
not addressed to the Society at all. On 7 July 2007 Benedict XVI issued
the motu proprio \work{Summorum Pontificum}, which---in the description
given by his own later act---broadened and adapted, ``by means of more
precisely determined norms,'' what \work{Ecclesia Dei} had contained in
general terms about the faculty of using the Roman Missal of
1962.\endnote{Benedict XVI, \work{Ecclesiae unitatem}, 2 July 2009, n.~3,
\work{Acta Apostolicae Sedis} 101 (2009) 710, with footnote 3 citing
\work{Summorum Pontificum} (7 July 2007), \work{Acta Apostolicae Sedis}
99 (2007) 777--781. The characterization quoted is the Holy See's own, in
its later act; the text of \work{Summorum Pontificum} itself is not
analyzed in this account, which is a history of the Society and not of
liturgical law. Its later abrogation and replacement by \work{Traditionis
custodes} (16 July 2021) lie outside this account's object and are noted
only because they form the background of the 2020s.}

That sequence---1984 indult, 2007 motu proprio---matters for one reason
here. Rome's liturgical concessions were real and were made without any
concession on status: throughout, the Society remained a body which, in
the words repeated in every Roman document from 1976 to 2026, has no
canonical status in the Church.
